\documentclass[10pt,a4paper]{amsart}
\usepackage[margin=19mm]{geometry}
\usepackage{amsmath,amssymb,mathtools}
\usepackage[scr=rsfs,cal=euler]{mathalfa}
\usepackage{xfrac}
\usepackage[unicode,hidelinks,bookmarks=false]{hyperref}
\newcommand{\ie}{i.e.\ }
\newcommand{\cf}{cf.\ }
\newcommand{\resp}{respectively\ }
\newcommand{\Subsection}{\subsection}
\providecommand{\nobreakdash}{\nobreak\hbox{-}}
\setlength{\emergencystretch}{2em}
\multlinegap0pt
\usepackage{amssymb}
\usepackage{algorithm}\usepackage{algpseudocode}
\usepackage{xcolor}
\newtheorem{proposition}{Proposition}
\newcommand{\QSent}{\mathbf{Q}} % Set of Q-sentences
\title{Goal-Oriented Logic-based Semantic Communication for Neuro-Symbolic Reasoning with Applications onto Autonomous Driving}
\author{Ahmet Faruk Saz, Duo Xu,, Faramarz Fekri}
\date{Source: arXiv:2608.00878v1 (CC BY 4.0)}
\begin{document}
\maketitle

\noindent\emph{Source and licence.} Goal-Oriented Logic-based Semantic Communication for Neuro-Symbolic Reasoning with Applications onto Autonomous Driving. Version arXiv:2608.00878v1 (CC BY 4.0), distributed by arXiv under the Creative Commons Attribution 4.0 International licence, http://creativecommons.org/licenses/by/4.0/, as recorded in the arXiv licence metadata for this exact version and shown on the abstract page https://arxiv.org/abs/2608.00878v1. Full licence notice: \texttt{licenses/arxiv-2608.00878-CC-BY-4.0.txt}.\par\smallskip

\noindent\textit{Editorial scope.} Counted unit: the article's proposition prop:q-partition (source0.tex lines 374--389). The article's complete original proof of it is delivered with it (source0.tex lines 391--396, one \texttt{proof} environment of 6 lines). No mathematics is written, repaired or re-derived: the statement and the proof are the article's own lines, in the article's own notation, and no retained line is substituted or re-flowed.  No further source-local context is needed: every cross-reference of every retained line resolves inside the two delivered spans.  Nothing was omitted from any retained span, so the package carries no omission record.  No retained line contains a citation, so the wrapper carries no bibliography and no bibliography entry.  No retained line contains a figure environment, an image-inclusion command or drawing code, so no image is delivered and the package declares no asset dependency.  Editorial formatting only, in addition: an AMS \texttt{amsart} preamble that restates the article's own declarations and macros used by the retained lines, namely the environments \texttt{proposition} on separate counters, together with the standard packages those lines need.  The retained environments print in this document as Proposition 1 in order of appearance, whereas the article prints the same items with the numbering of its own full document.  The complete native source of the article as distributed in the arXiv e-print for this exact version is bundled unchanged as \texttt{sources/206573/source0.tex}.
\begin{proposition}[Q-Sentence Partition]
\label{prop:q-partition}
For any assignment of two entities (i.e., individuals) $v_1$ and $v_2$ to variables $x_1$ and $x_2$, exactly one Q-sentence in $\QSent$ is satisfied.
Equivalently,
\begin{align}
    \bigvee_{i=1}^{K}Q_i(v_1,v_2)
    &\equiv \top,
    \label{eq:q-exhaustive}\\
    Q_i(v_1,v_2)\land Q_j(v_1,v_2)
    &\equiv \bot,
    \qquad i\neq j.
    \label{eq:q-exclusive}
\end{align}

That particular assignment of entities to variables instantiates that Q-sentence.
\end{proposition}

\begin{proof}
Under a fixed model and valuation, each predicate slot $A_t$ is either true
or false. These truth values determine a unique binary vector in
$\{0,1\}^{T}$ and hence a unique Q-sentence. Any two distinct Q-sentences
differ in at least one predicate slot and therefore cannot both be true.
\end{proof}

\end{document}
